% Options for packages loaded elsewhere
\PassOptionsToPackage{unicode}{hyperref}
\PassOptionsToPackage{hyphens}{url}
\PassOptionsToPackage{dvipsnames,svgnames,x11names}{xcolor}
\documentclass[
  conference]{IEEEtran}
\usepackage{xcolor}
\usepackage{amsmath,amssymb}
\setcounter{secnumdepth}{-\maxdimen} % remove section numbering
\usepackage{iftex}
\ifPDFTeX
  \usepackage[T1]{fontenc}
  \usepackage[utf8]{inputenc}
  \usepackage{textcomp} % provide euro and other symbols
\else % if luatex or xetex
  \usepackage{unicode-math} % this also loads fontspec
  \defaultfontfeatures{Scale=MatchLowercase}
  \defaultfontfeatures[\rmfamily]{Ligatures=TeX,Scale=1}
\fi
\usepackage{lmodern}
\ifPDFTeX\else
  % xetex/luatex font selection
\fi
% Use upquote if available, for straight quotes in verbatim environments
\IfFileExists{upquote.sty}{\usepackage{upquote}}{}
\IfFileExists{microtype.sty}{% use microtype if available
  \usepackage[]{microtype}
  \UseMicrotypeSet[protrusion]{basicmath} % disable protrusion for tt fonts
}{}
\makeatletter
\@ifundefined{KOMAClassName}{% if non-KOMA class
  \IfFileExists{parskip.sty}{%
    \usepackage{parskip}
  }{% else
    \setlength{\parindent}{0pt}
    \setlength{\parskip}{6pt plus 2pt minus 1pt}}
}{% if KOMA class
  \KOMAoptions{parskip=half}}
\makeatother
\usepackage{longtable,booktabs,array}
\usepackage{caption}
\captionsetup[table]{skip=6pt}
\newcounter{none} % for unnumbered tables
\usepackage{calc} % for calculating minipage widths
% Correct order of tables after \paragraph or \subparagraph
\usepackage{etoolbox}
\makeatletter
\patchcmd\longtable{\par}{\if@noskipsec\mbox{}\fi\par}{}{}
\makeatother
% definitions for citeproc citations
\NewDocumentCommand\citeproctext{}{}
\NewDocumentCommand\citeproc{mm}{%
  \begingroup\def\citeproctext{#2}\cite{#1}\endgroup}
\makeatletter
 % allow citations to break across lines
 \let\@cite@ofmt\@firstofone
 % avoid brackets around text for \cite:
 \def\@biblabel#1{}
 \def\@cite#1#2{{#1\if@tempswa , #2\fi}}
\makeatother
\newlength{\cslhangindent}
\setlength{\cslhangindent}{1.5em}
\newlength{\csllabelwidth}
\setlength{\csllabelwidth}{3em}
\newenvironment{CSLReferences}[2] % #1 hanging-indent, #2 entry-spacing
 {\begin{list}{}{%
  \setlength{\itemindent}{0pt}
  \setlength{\leftmargin}{0pt}
  \setlength{\parsep}{0pt}
  % turn on hanging indent if param 1 is 1
  \ifodd #1
   \setlength{\leftmargin}{\cslhangindent}
   \setlength{\itemindent}{-1\cslhangindent}
  \fi
  % set entry spacing
  \setlength{\itemsep}{#2\baselineskip}}}
 {\end{list}}
\usepackage{calc}
\newcommand{\CSLBlock}[1]{\hfill\break\parbox[t]{\linewidth}{\strut\ignorespaces#1\strut}}
\newcommand{\CSLLeftMargin}[1]{\parbox[t]{\csllabelwidth}{\strut#1\strut}}
\newcommand{\CSLRightInline}[1]{\parbox[t]{\linewidth - \csllabelwidth}{\strut#1\strut}}
\newcommand{\CSLIndent}[1]{\hspace{\cslhangindent}#1}
\setlength{\emergencystretch}{3em} % prevent overfull lines
\providecommand{\tightlist}{%
  \setlength{\itemsep}{0pt}\setlength{\parskip}{0pt}}
\usepackage{times}
\usepackage{bookmark}
\IfFileExists{xurl.sty}{\usepackage{xurl}}{} % add URL line breaks if available
\urlstyle{same}
% fallback for those not using the hyperref driver hyperxmp:
\makeatletter
\@ifundefined{xmpquote}{\newcommand{\xmpquote}[1]{#1}}{}
\makeatother
\hypersetup{
  pdftitle={Stability-Penalized Split Selection for Decision Tree Classification: Sebuah Evaluasi Multi-Dataset yang Reproducible},
  pdfauthor={Rofiq Fauzi (p36202600047@mhs.dinus.ac.id), Magister Artificial Intelligence, Universitas Dian Nuswantoro},
  colorlinks=true,
  linkcolor={black},
  filecolor={Maroon},
  citecolor={Blue},
  urlcolor={black},
  pdfcreator={LaTeX via pandoc}}

\title{Stability-Penalized Split Selection for Decision Tree
Classification: Sebuah Evaluasi Multi-Dataset yang Reproducible}
\author{Rofiq Fauzi (p36202600047@mhs.dinus.ac.id), Magister Artificial
Intelligence, Universitas Dian Nuswantoro}
\date{}

\begin{document}
\maketitle

\textbf{Kata kunci}--pohon keputusan; kriteria split; stabilitas
bootstrap; variansi gain; \emph{nested cross-validation}; perbandingan
statistik nonparametrik.

\subsection{Abstrak}\label{abstrak}

Pohon keputusan CART dikenal cepat dan interpretabel, namun struktur
pohonnya tidak stabil: perubahan kecil pada data latih dapat mengubah
split yang terpilih dan akhirnya prediksi. Penelitian ini mengevaluasi
sebuah formulasi kandidat yang sangat sederhana, dinamai
\emph{Stability-Penalized Decision Tree} (SPDT), di mana skor setiap
kandidat split dihitung sebagai gain Gini dikurangi penalti berbobot
terhadap variansi gain tersebut atas resample bootstrap dari data node:
\texttt{Score(s)\ =\ Gain(s,D)\ -\ lambda*Var\_b{[}Gain(s,D\^{}(b)){]}}.
Penelusuran literatur hingga 22 Oktober 2026 tidak menemukan ekuivalen
formula-level dari kriteria tersebut, tetapi ruang riset stabilitas
pohon sudah ramai; karena itu formulasi diposisikan sebagai kandidat,
bukan klaim metode baru. Evaluasi dilakukan pada 8 dataset klasifikasi
tabular UCI dengan \emph{nested cross-validation} berulang (5 fold x 3
ulangan pada tiga seed), penyetelan hyperparameter di inner loop, dan
Macro-F1 sebagai metrik primer. Hasil: SPDT vs CART custom (lambda=0)
mencatat 6 menang, 0 seri, 2 kalah dengan rata-rata Delta Macro-F1
+0.0058. Uji Friedman menunjukkan chi\^{}2(2)=4.0000, p=0.1353, Kendall
W=0.2500 -- tidak cukup bukti untuk menolak hipotesis nol, sehingga
perbandingan post-hoc konfirmatori tidak dijalankan. Ablation
menunjukkan besar lambda dan jumlah resample B nyaris tidak memengaruhi
Macro-F1, pembatasan kedalaman lebih menentukan, dan ukuran stabilitas
prediksi SPDT sedikit lebih buruk daripada CART custom dengan overhead
waktu fit +/-40x. Temuan ini merupakan hasil nol yang jujur: pada
delapan dataset kecil-menengah yang diuji, penalti variansi bootstrap
atas gain split tidak memberikan manfaat yang cukup untuk menandingi
kesederhanaan CART, tanpa menutup kemungkinan manfaat pada regime data
lain.

\subsection{I. Pendahuluan}\label{i.-pendahuluan}

Pohon keputusan adalah metode klasifikasi nonparametrik yang mempartisi
ruang fitur secara rekursif. Kriteria split CART memilih kandidat
(fitur, ambang) yang memaksimalkan penurunan impuritas Gini (Breiman et
al. 1984; Quinlan 1986). Namun demikian, struktur pohon diketahui tidak
stabil: perturbasi kecil pada data latih dapat mengubah split yang
terpilih (Li and Clarke 2002), dan konsekuensinya prediksi berubah --
fenomena inilah yang melatarbelakangi \emph{bagging} dan \emph{random
forests} (Breiman 1996, 2001; Dietterich 2000).

Berbagai pendekatan telah mengurangi instabilitas tersebut pada tingkat
yang berbeda: filter signifikansi statistik, pembatasan struktural,
agregasi ensemble, atau pemilihan model yang stabil (Briand et al. 2009;
Lee et al. 2024; Mirzamomen and Kangavari 2017; Wang et al. 2017). Yang
jarang diuji secara langsung adalah pertanyaan yang lebih sempit ini:

\begin{quote}
\textbf{Apakah penalti variansi bootstrap atas gain split, dikurangkan
langsung dari skor split, dapat meningkatkan Macro-F1 out-of-sample
dan/atau stabilitas prediksi dibandingkan CART standar pada resampling
berulang, dengan biaya runtime yang dapat diterima?}
\end{quote}

Hipotesis nol: tidak ada peningkatan sistematis dibandingkan baseline
lintas dataset independen. Tidak ada hipotesis bahwa SPDT \emph{harus}
unggul.

Kontribusi yang diklaim terbatas pada: (i) formulasi kandidat yang
eksplisit; (ii) implementasi dan audit reproducible; (iii) evaluasi
nested-CV yang bebas \emph{leakage}; (iv) ablation dan analisis
statistik yang jujur terhadap hasil positif maupun negatif.

\subsection{II. Karya Terkait}\label{ii.-karya-terkait}

\textbf{Stabilitas pohon keputusan.} Instabilitas struktural pohon
dibahas secara formal oleh Li dan Clarke (Li and Clarke 2002);
pengukuran kuantitatif stabilitas dan kemiripan antar pohon dikembangkan
antara lain oleh Turney (Turney 1995) serta Briand dkk. (Briand et al.
2009), yang juga mengevaluasi stabilisasi node berbasis bootstrap (lini
NLS/Dannegger). Lee dkk. (Lee et al. 2024) mengusulkan ANLS yang
menjalankan iterasi bootstrap pada setiap node dan memilih prediktor
terpilih terbanyak -- mekanisme resampling-node yang paling dekat dengan
penelitian ini, tetapi agregasinya memakai voting frekuensi, bukan
penalti variansi gain. Mirzamomen dan Kangavari (Mirzamomen and
Kangavari 2017) menstabilkan melalui uji split multi-atribut.

\textbf{Kriteria split alternatif.} Kriteria berbasis entropi Tsallis
dua-suku (Wang et al. 2017), kriteria \emph{robust} terhadap label noise
(Yang et al. 2019), pemilihan split tak-bias berbasis statistik Gini
(Strobl et al. 2007), dan pohon arus data dengan batas keyakinan
Hoeffding (Domingos and Hulten 2000) menunjukkan bahwa ruang ``mengubah
skor split'' sudah aktif diteliti. Delise (Delise 2026) menilai
konsistensi gain lintas \emph{environment} pada boosting. Penalti skor
oleh variansi estimasi resampled juga dikenal di luar pohon (pola
mean-variance pada seleksi model). Tidak satu pun dari karya yang
berhasil diverifikasi memuat rumus
\texttt{Gain\ -\ lambda*Var\_bootstrap{[}Gain{]}}; namun demikian temuan
negatif penelusuran bukan bukti kebaruan, sehingga formulasi tetap dicap
sebagai kandidat.

\subsection{III. Metode}\label{iii.-metode}

\subsubsection{A. Kriteria split yang
diusulkan}\label{a.-kriteria-split-yang-diusulkan}

Untuk node dengan data D (hanya bagian latih yang sah), dengan
\texttt{Gini(D)\ =\ 1\ -\ Sum\_c\ p(c\textbar{}D)\^{}2} dan gain split

\texttt{Gain(s,D)\ =\ Gini(D)\ -\ (\textbar{}D\_L\textbar{}/\textbar{}D\textbar{})*Gini(D\_L)\ -\ (\textbar{}D\_R\textbar{}/\textbar{}D\textbar{})*Gini(D\_R)},

dibangun B resample bootstrap \texttt{D\^{}(b)} dari baris node. Untuk
setiap kandidat s yang valid:

\texttt{V(s)\ =\ Var\_b{[}Gain(s,D\^{}(b)){]}} (varians sampel, ddof=1)

\textbf{Skor utama: \texttt{Score(s)\ =\ Gain(s,D)\ -\ lambda*V(s)}}.

Kebijakan deterministik: ambang kandidat = titik tengah nilai unik
terurut; \texttt{min\_samples\_leaf}/\texttt{min\_samples\_split} pada
jumlah asli node; partisi child kosong dalam resample berkontribusi 0;
daun memilih kelas mayoritas dengan tie-break label terkecil; seed
bootstrap diturunkan dari \texttt{(random\_state,\ id\ node)}. Untuk
lambda=0 seluruh komputasi bootstrap dilewati dan estimator
terdegenerasi menjadi CART custom referensi (diverifikasi identik dengan
implementasi referensi independen pada pengujian unit). Sebuah varian
eksploratori terpisah (normalisasi \texttt{V/(Gain\^{}2+eps)}) dicatat
tetapi bukan metode utama. Keterbatasan penting: V(s)
\emph{scale-dependent} terhadap gain.

\subsubsection{B. Desain evaluasi}\label{b.-desain-evaluasi}

Dataset: 8 dataset klasifikasi tabular UCI (Tabel I). Protokol:
\emph{repeated stratified nested CV} -- outer 5 fold x 3 ulangan untuk
masing-masing seed {[}42, 123, 2026{]} (45 evaluasi per dataset-metode,
split identik untuk semua metode); inner 3-fold stratified untuk
penyetelan dengan seed deterministik. Preprocessing hanya imputasi
median yang di-fit pada bagian latih (pipeline). Metrik primer Macro-F1;
sekunder: balanced accuracy, accuracy, waktu fit/prediksi, jumlah node,
kedalaman. Stabilitas prediksi diukur terpisah sebagai
\emph{disagreement} rata-rata pairwise antar model yang dilatih pada 10
sub-sampel 90\% dari data latih outer, dievaluasi pada outer test yang
sama -- dilaporkan terpisah dari akurasi.

\begin{longtable}[]{@{}lllll@{}}
\caption{Inventaris dataset.}\tabularnewline
\toprule\noalign{}
Dataset & N & Fitur & Kelas & Sumber \\
\midrule\noalign{}
\endfirsthead
\toprule\noalign{}
Dataset & N & Fitur & Kelas & Sumber \\
\midrule\noalign{}
\endhead
\bottomrule\noalign{}
\endlastfoot
banknote & 1372 & 4 & 2 & ucimlrepo \\
glass & 214 & 9 & 6 & ucimlrepo \\
ionosphere & 351 & 34 & 2 & ucimlrepo \\
iris & 150 & 4 & 3 & ucimlrepo \\
seeds & 210 & 7 & 3 & uci\_static\_zip \\
sonar & 208 & 60 & 2 & ucimlrepo \\
wdbc & 569 & 30 & 2 & ucimlrepo \\
wine & 178 & 13 & 3 & ucimlrepo \\
\end{longtable}

\subsubsection{C. Analisis statistik}\label{c.-analisis-statistik}

Unit inferensi adalah dataset (rata-rata 45 outer fold), bukan fold.
Perbandingan tiga metode primer (CART sklearn, CART custom lambda=0,
SPDT) menggunakan uji Friedman dengan dataset sebagai blok; bila omnibus
memberikan alasan, dilanjutkan Wilcoxon signed-rank berpasangan dengan
koreksi Holm dan effect size rank-biserial. p \textgreater= 0.05 tidak
pernah diartikan sebagai bukti kesetaraan. Random forest disertakan
sebagai benchmark kontekstual non-ekuivalen dan tidak masuk uji omnibus
primer.

\subsection{IV. Hasil dan Pembahasan}\label{iv.-hasil-dan-pembahasan}

\begin{longtable}[]{@{}
  >{\raggedright\arraybackslash}p{(\linewidth - 10\tabcolsep) * \real{0.1667}}
  >{\raggedright\arraybackslash}p{(\linewidth - 10\tabcolsep) * \real{0.1667}}
  >{\raggedright\arraybackslash}p{(\linewidth - 10\tabcolsep) * \real{0.1667}}
  >{\raggedright\arraybackslash}p{(\linewidth - 10\tabcolsep) * \real{0.1667}}
  >{\raggedright\arraybackslash}p{(\linewidth - 10\tabcolsep) * \real{0.1667}}
  >{\raggedright\arraybackslash}p{(\linewidth - 10\tabcolsep) * \real{0.1667}}@{}}
\caption{Macro-F1 mean +/- sd (deskriptif, fold-level) per dataset dan
metode.}\tabularnewline
\toprule\noalign{}
\begin{minipage}[b]{\linewidth}\raggedright
Dataset
\end{minipage} & \begin{minipage}[b]{\linewidth}\raggedright
cart\_sklearn
\end{minipage} & \begin{minipage}[b]{\linewidth}\raggedright
cart\_custom
\end{minipage} & \begin{minipage}[b]{\linewidth}\raggedright
spdt
\end{minipage} & \begin{minipage}[b]{\linewidth}\raggedright
random\_forest
\end{minipage} & \begin{minipage}[b]{\linewidth}\raggedright
Delta (spdt-cart\_custom)
\end{minipage} \\
\midrule\noalign{}
\endfirsthead
\toprule\noalign{}
\begin{minipage}[b]{\linewidth}\raggedright
Dataset
\end{minipage} & \begin{minipage}[b]{\linewidth}\raggedright
cart\_sklearn
\end{minipage} & \begin{minipage}[b]{\linewidth}\raggedright
cart\_custom
\end{minipage} & \begin{minipage}[b]{\linewidth}\raggedright
spdt
\end{minipage} & \begin{minipage}[b]{\linewidth}\raggedright
random\_forest
\end{minipage} & \begin{minipage}[b]{\linewidth}\raggedright
Delta (spdt-cart\_custom)
\end{minipage} \\
\midrule\noalign{}
\endhead
\bottomrule\noalign{}
\endlastfoot
banknote & 0.982 +/- 0.008 & 0.981 +/- 0.007 & 0.981 +/- 0.007 & 0.994
+/- 0.004 & +0.0004 \\
glass & 0.620 +/- 0.097 & 0.609 +/- 0.093 & 0.626 +/- 0.093 & 0.742 +/-
0.075 & +0.0166 \\
ionosphere & 0.873 +/- 0.034 & 0.881 +/- 0.038 & 0.879 +/- 0.038 & 0.929
+/- 0.037 & -0.0027 \\
iris & 0.949 +/- 0.033 & 0.945 +/- 0.038 & 0.944 +/- 0.037 & 0.955 +/-
0.034 & -0.0008 \\
seeds & 0.897 +/- 0.032 & 0.906 +/- 0.036 & 0.913 +/- 0.034 & 0.923 +/-
0.038 & +0.0066 \\
sonar & 0.719 +/- 0.065 & 0.706 +/- 0.073 & 0.725 +/- 0.061 & 0.833 +/-
0.065 & +0.0189 \\
wdbc & 0.925 +/- 0.031 & 0.925 +/- 0.030 & 0.928 +/- 0.030 & 0.956 +/-
0.023 & +0.0037 \\
wine & 0.905 +/- 0.047 & 0.905 +/- 0.044 & 0.909 +/- 0.046 & 0.981 +/-
0.024 & +0.0040 \\
\end{longtable}

\begin{longtable}[]{@{}ll@{}}
\caption{Mean rank Macro-F1.}\tabularnewline
\toprule\noalign{}
Metode & Mean rank \\
\midrule\noalign{}
\endfirsthead
\toprule\noalign{}
Metode & Mean rank \\
\midrule\noalign{}
\endhead
\bottomrule\noalign{}
\endlastfoot
spdt & 1.5000 \\
cart\_sklearn & 2.0000 \\
cart\_custom & 2.5000 \\
\end{longtable}

\begin{longtable}[]{@{}
  >{\raggedright\arraybackslash}p{(\linewidth - 12\tabcolsep) * \real{0.1429}}
  >{\raggedright\arraybackslash}p{(\linewidth - 12\tabcolsep) * \real{0.1429}}
  >{\raggedright\arraybackslash}p{(\linewidth - 12\tabcolsep) * \real{0.1429}}
  >{\raggedright\arraybackslash}p{(\linewidth - 12\tabcolsep) * \real{0.1429}}
  >{\raggedright\arraybackslash}p{(\linewidth - 12\tabcolsep) * \real{0.1429}}
  >{\raggedright\arraybackslash}p{(\linewidth - 12\tabcolsep) * \real{0.1429}}
  >{\raggedright\arraybackslash}p{(\linewidth - 12\tabcolsep) * \real{0.1429}}@{}}
\caption{Uji statistik dataset-level.}\tabularnewline
\toprule\noalign{}
\begin{minipage}[b]{\linewidth}\raggedright
Uji
\end{minipage} & \begin{minipage}[b]{\linewidth}\raggedright
Statistik
\end{minipage} & \begin{minipage}[b]{\linewidth}\raggedright
db
\end{minipage} & \begin{minipage}[b]{\linewidth}\raggedright
p
\end{minipage} & \begin{minipage}[b]{\linewidth}\raggedright
\#dataset
\end{minipage} & \begin{minipage}[b]{\linewidth}\raggedright
Kendall W
\end{minipage} & \begin{minipage}[b]{\linewidth}\raggedright
Keputusan
\end{minipage} \\
\midrule\noalign{}
\endfirsthead
\toprule\noalign{}
\begin{minipage}[b]{\linewidth}\raggedright
Uji
\end{minipage} & \begin{minipage}[b]{\linewidth}\raggedright
Statistik
\end{minipage} & \begin{minipage}[b]{\linewidth}\raggedright
db
\end{minipage} & \begin{minipage}[b]{\linewidth}\raggedright
p
\end{minipage} & \begin{minipage}[b]{\linewidth}\raggedright
\#dataset
\end{minipage} & \begin{minipage}[b]{\linewidth}\raggedright
Kendall W
\end{minipage} & \begin{minipage}[b]{\linewidth}\raggedright
Keputusan
\end{minipage} \\
\midrule\noalign{}
\endhead
\bottomrule\noalign{}
\endlastfoot
friedman & 4.0000 & 2 & 0.1353 & 8 & 0.2500 & tidak ada bukti omnibus;
post-hoc tidak dijalankan \\
\end{longtable}

\emph{Post-hoc tidak dijalankan: uji omnibus Friedman tidak memberikan
alasan (p \textgreater= 0.05), sesuai protokol.}

Table: Perbandingan berpasangan (post-hoc).

Ringkasan berpasangan SPDT vs CART custom: \textbf{6 menang / 0 seri / 2
kalah}, rata-rata Delta Macro-F1 +0.0058.

\begin{longtable}[]{@{}
  >{\raggedright\arraybackslash}p{(\linewidth - 10\tabcolsep) * \real{0.1667}}
  >{\raggedright\arraybackslash}p{(\linewidth - 10\tabcolsep) * \real{0.1667}}
  >{\raggedright\arraybackslash}p{(\linewidth - 10\tabcolsep) * \real{0.1667}}
  >{\raggedright\arraybackslash}p{(\linewidth - 10\tabcolsep) * \real{0.1667}}
  >{\raggedright\arraybackslash}p{(\linewidth - 10\tabcolsep) * \real{0.1667}}
  >{\raggedright\arraybackslash}p{(\linewidth - 10\tabcolsep) * \real{0.1667}}@{}}
\caption{Biaya komputasi.}\tabularnewline
\toprule\noalign{}
\begin{minipage}[b]{\linewidth}\raggedright
Metode
\end{minipage} & \begin{minipage}[b]{\linewidth}\raggedright
Fit rata-rata (s)
\end{minipage} & \begin{minipage}[b]{\linewidth}\raggedright
Fit median (s)
\end{minipage} & \begin{minipage}[b]{\linewidth}\raggedright
Prediksi rata-rata (s)
\end{minipage} & \begin{minipage}[b]{\linewidth}\raggedright
Node rata-rata
\end{minipage} & \begin{minipage}[b]{\linewidth}\raggedright
Kedalaman rata-rata
\end{minipage} \\
\midrule\noalign{}
\endfirsthead
\toprule\noalign{}
\begin{minipage}[b]{\linewidth}\raggedright
Metode
\end{minipage} & \begin{minipage}[b]{\linewidth}\raggedright
Fit rata-rata (s)
\end{minipage} & \begin{minipage}[b]{\linewidth}\raggedright
Fit median (s)
\end{minipage} & \begin{minipage}[b]{\linewidth}\raggedright
Prediksi rata-rata (s)
\end{minipage} & \begin{minipage}[b]{\linewidth}\raggedright
Node rata-rata
\end{minipage} & \begin{minipage}[b]{\linewidth}\raggedright
Kedalaman rata-rata
\end{minipage} \\
\midrule\noalign{}
\endhead
\bottomrule\noalign{}
\endlastfoot
cart\_custom & 0.5340 & 0.4163 & 0.0009 & 26.6722 & 5.1000 \\
cart\_sklearn & 0.2693 & 0.2605 & 0.0010 & 27.4222 & 5.2917 \\
random\_forest & 1.1169 & 1.0863 & 0.0573 & nan & nan \\
spdt & 21.4302 & 14.6491 & 0.0010 & 28.5944 & 5.4611 \\
\end{longtable}

\begin{longtable}[]{@{}lll@{}}
\caption{Stabilitas prediksi (disagreement; bukan metrik
akurasi).}\tabularnewline
\toprule\noalign{}
Metode & Disagreement rata-rata & SD \\
\midrule\noalign{}
\endfirsthead
\toprule\noalign{}
Metode & Disagreement rata-rata & SD \\
\midrule\noalign{}
\endhead
\bottomrule\noalign{}
\endlastfoot
cart\_custom & 0.0928 & 0.0853 \\
cart\_sklearn & 0.0983 & 0.0885 \\
random\_forest & 0.0271 & 0.0337 \\
spdt & 0.1014 & 0.0934 \\
\end{longtable}

\begin{longtable}[]{@{}
  >{\raggedright\arraybackslash}p{(\linewidth - 8\tabcolsep) * \real{0.2000}}
  >{\raggedright\arraybackslash}p{(\linewidth - 8\tabcolsep) * \real{0.2000}}
  >{\raggedright\arraybackslash}p{(\linewidth - 8\tabcolsep) * \real{0.2000}}
  >{\raggedright\arraybackslash}p{(\linewidth - 8\tabcolsep) * \real{0.2000}}
  >{\raggedright\arraybackslash}p{(\linewidth - 8\tabcolsep) * \real{0.2000}}@{}}
\caption{Ablation pada split identik.}\tabularnewline
\toprule\noalign{}
\begin{minipage}[b]{\linewidth}\raggedright
Varian
\end{minipage} & \begin{minipage}[b]{\linewidth}\raggedright
F1 makro rata-rata
\end{minipage} & \begin{minipage}[b]{\linewidth}\raggedright
SD
\end{minipage} & \begin{minipage}[b]{\linewidth}\raggedright
Fit rata-rata (s)
\end{minipage} & \begin{minipage}[b]{\linewidth}\raggedright
Node rata-rata
\end{minipage} \\
\midrule\noalign{}
\endfirsthead
\toprule\noalign{}
\begin{minipage}[b]{\linewidth}\raggedright
Varian
\end{minipage} & \begin{minipage}[b]{\linewidth}\raggedright
F1 makro rata-rata
\end{minipage} & \begin{minipage}[b]{\linewidth}\raggedright
SD
\end{minipage} & \begin{minipage}[b]{\linewidth}\raggedright
Fit rata-rata (s)
\end{minipage} & \begin{minipage}[b]{\linewidth}\raggedright
Node rata-rata
\end{minipage} \\
\midrule\noalign{}
\endhead
\bottomrule\noalign{}
\endlastfoot
A0\_lam0 & 0.8599 & 0.1270 & 0.0201 & 37.6611 \\
A1\_lam0.1\_B10 & 0.8598 & 0.1256 & 0.0860 & 37.6111 \\
A1\_lam1\_B10 & 0.8601 & 0.1255 & 0.0862 & 37.5278 \\
A2\_B20 & 0.8601 & 0.1251 & 0.1136 & 37.4556 \\
A2\_B5 & 0.8613 & 0.1243 & 0.0711 & 37.4667 \\
A3\_d3\_leaf1 & 0.8357 & 0.1592 & 0.0371 & 12.8722 \\
A3\_d3\_leaf5 & 0.8353 & 0.1546 & 0.0373 & 11.8611 \\
A3\_d8\_leaf1 & 0.8592 & 0.1272 & 0.0821 & 35.2500 \\
A3\_d8\_leaf5 & 0.8521 & 0.1381 & 0.0820 & 24.7389 \\
\end{longtable}

\begin{longtable}[]{@{}llll@{}}
\caption{Frekuensi seleksi hyperparameter SPDT (45 outer fold per
dataset).}\tabularnewline
\toprule\noalign{}
parameter & value & frequency & share \\
\midrule\noalign{}
\endfirsthead
\toprule\noalign{}
parameter & value & frequency & share \\
\midrule\noalign{}
\endhead
\bottomrule\noalign{}
\endlastfoot
max\_depth & 5.0000 & 140 & 0.3889 \\
max\_depth & 8.0000 & 102 & 0.2833 \\
max\_depth & 3.0000 & 92 & 0.2556 \\
min\_samples\_leaf & 1.0000 & 201 & 0.5583 \\
min\_samples\_leaf & 2.0000 & 97 & 0.2694 \\
min\_samples\_leaf & 5.0000 & 62 & 0.1722 \\
min\_samples\_split & 2.0000 & 252 & 0.7000 \\
min\_samples\_split & 5.0000 & 69 & 0.1917 \\
min\_samples\_split & 10.0000 & 39 & 0.1083 \\
n\_bootstrap & 5.0000 & 191 & 0.5306 \\
n\_bootstrap & 20.0000 & 87 & 0.2417 \\
n\_bootstrap & 10.0000 & 82 & 0.2278 \\
penalty\_lambda & 10.0000 & 233 & 0.6472 \\
penalty\_lambda & 0.0000 & 61 & 0.1694 \\
penalty\_lambda & 0.0100 & 41 & 0.1139 \\
penalty\_lambda & 1.0000 & 22 & 0.0611 \\
penalty\_lambda & 0.1000 & 3 & 0.0083 \\
\end{longtable}

\textbf{Pembahasan.} Hasil bersifat deskriptif condong mendukung SPDT
tetapi tidak mencapai bukti konfirmatori. SPDT mencatat mean rank
terbaik (1.500 vs 2.000 CART sklearn dan 2.500 CART custom) dan menang
pada 6 dari 8 dataset berpasangan terhadap CART custom, namun uji
Friedman tidak menemukan bukti perbedaan yang cukup (p = 0.1353) dan
perbandingan eksploratif SPDT-CART custom menghasilkan p mentah 0.078 --
tidak bertahan pada taraf 5\% dan tidak boleh ditafsirkan sebagai
signifikan. Dengan hanya 8 dataset, power statistik terbatas; sesuai
protokol, hasil ini dilaporkan apa adanya dan p \textgreater= 0.05 tidak
diartikan sebagai kesetaraan.

Temuan penting kedua: penalti stabilitas tidak menurunkan
\emph{prediction disagreement} -- SPDT justru sedikit lebih tinggi
(0.101) dibanding CART custom (0.093) dan CART sklearn (0.101). Random
forest tetap paling stabil (0.027) melalui agregasi ensemble. Dengan
demikian, janji motivasi metode (stabilitas prediksi lebih baik)
\textbf{tidak terwujud} pada benchmark ini; penalti variansi gain bukan
pengganti ensemble.

Temuan ketiga dari ablation pada split identik: besar penalti lambda
maupun jumlah resample B nyaris tidak mengubah Macro-F1 (semua varian
lambda dan B di kisaran 0.859-0.861), sedangkan pembatasan
kedalaman/ukuran daun jauh lebih menentukan (depth 3 turun ke +/-0.836).
Penyetelan inner cenderung memilih lambda besar (10 pada 64.7\% fold)
dan B kecil (5 pada 53.1\%), yang mengindikasikan penalti sering dipilih
tetapi efeknya pada akurasi kecil. Biayanya nyata: waktu fit SPDT +/-40x
CART custom. Secara keseluruhan, efektivitas penalti variansi bootstrap
bergantung dataset dan kecil magnitude-nya; simple distance/impurity
splitting tetap preferable untuk dataset kecil-menengah ini.

\subsection{V. Kesimpulan}\label{v.-kesimpulan}

Penelitian ini mengevaluasi formulasi kandidat SPDT secara reproducible
dengan nested repeated CV bebas leakage, ablation terkontrol pada split
identik, dan analisis statistik level-dataset pada 8 dataset UCI. Secara
deskriptif SPDT memperoleh mean rank terbaik dan menang berpasangan pada
6 dari 8 dataset terhadap CART custom, tetapi uji Friedman tidak
memberikan bukti perbedaan yang cukup (chi\^{}2(2)=4.0000, p=0.1353),
sehingga tidak ada dasar konfirmatori untuk mengklaim peningkatan. Lebih
lanjut, ukuran stabilitas prediksi tidak membaik dan justru sedikit
memburuk, sementara biaya komputasi naik +/-40x. Pada kondisi yang
diuji, kesederhanaan CART tetap preferable. Keterbatasan: delapan
dataset kecil, impuritas Gini saja, proxy stabilitas berbasis
disagreement, klaim kebaruan formula-level yang belum established, dan
belum dibandingkannya kriteria terdekat dari literatur (ANLS, minimax
split). Pekerjaan lanjutan yang paling bernilai: membandingkan langsung
dengan stabilisasi bootstrap-node dari literatur, mengkaji normalisasi
penalti oleh skala gain, dan menguji pada benchmark yang lebih besar.

\subsection{Referensi}\label{referensi}

\protect\phantomsection\label{refs}
\begin{CSLReferences}{1}{1}
\bibitem[\citeproctext]{ref-breiman1996bagging}
Breiman, Leo. 1996. {``Bagging Predictors.''} \emph{Machine Learning} 24
(2): 123--40. \url{https://doi.org/10.1023/A:1018054314350}.

\bibitem[\citeproctext]{ref-breiman2001randomforests}
Breiman, Leo. 2001. {``Random Forests.''} \emph{Machine Learning} 45
(1): 5--32. \url{https://doi.org/10.1023/A:1010933404324}.

\bibitem[\citeproctext]{ref-breiman1984cart}
Breiman, Leo, Jerome H. Friedman, Richard A. Olshen, and Charles J.
Stone. 1984. \emph{Classification and Regression Trees}. Wadsworth;
expanded edition Chapman \& Hall/CRC, 1993.

\bibitem[\citeproctext]{ref-briand2009stability}
Briand, Bénédicte, Gilles R. Ducharme, Vanessa Parache, and Catherine
Mercat-Rommens. 2009. {``A Similarity Measure to Assess the Stability of
Classification Trees.''} \emph{Computational Statistics \& Data
Analysis} 53 (4): 1208--17.
\url{https://doi.org/10.1016/j.csda.2008.10.033}.

\bibitem[\citeproctext]{ref-delise2026era}
Delise, Timothy. 2026. {``Era Splitting -- Invariant Learning for
Decision Trees.''} \emph{Expert Systems with Applications} 299: 129687.
\url{https://doi.org/10.1016/j.eswa.2025.129687}.

\bibitem[\citeproctext]{ref-dietterich2000ensemble}
Dietterich, Thomas G. 2000. {``An Experimental Comparison of Three
Methods for Constructing Ensembles of Decision Trees: Bagging, Boosting,
and Randomization.''} \emph{Machine Learning} 40 (2): 139--57.
\url{https://doi.org/10.1023/A:1007607513941}.

\bibitem[\citeproctext]{ref-domingos2000vfdt}
Domingos, Pedro, and Geoff Hulten. 2000. {``Mining High-Speed Data
Streams.''} \emph{Proceedings of the Sixth ACM SIGKDD International
Conference on Knowledge Discovery and Data Mining (KDD '00)}, 71--80.
\url{https://doi.org/10.1145/347090.347107}.

\bibitem[\citeproctext]{ref-lee2024anls}
Lee, Joon, Min Kyu Sim, and Jin-Seon Hong. 2024. {``Assessing Decision
Tree Stability: A Comprehensive Method for Generating a Stable Decision
Tree.''} \emph{IEEE Access} 12: 90061--72.
\url{https://doi.org/10.1109/ACCESS.2024.3419228}.

\bibitem[\citeproctext]{ref-li2002instability}
Li, Rong-Hua, and Stewart Clarke. 2002. {``Instability of Decision Tree
Classification Algorithms.''} \emph{Proceedings of the Eighth ACM SIGKDD
International Conference on Knowledge Discovery and Data Mining (KDD
'02)}, 570--75. \url{https://doi.org/10.1145/775047.775131}.

\bibitem[\citeproctext]{ref-mirzamomen2017stable}
Mirzamomen, Zahra, and Mohammad Reza Kangavari. 2017. {``A Framework to
Induce More Stable Decision Trees for Pattern Classification.''}
\emph{Pattern Analysis and Applications} 20 (4): 991--1004.
\url{https://doi.org/10.1007/s10044-016-0542-2}.

\bibitem[\citeproctext]{ref-quinlan1986induction}
Quinlan, J. Ross. 1986. {``Induction of Decision Trees.''} \emph{Machine
Learning} 1 (1): 81--106. \url{https://doi.org/10.1023/A:1022643204877}.

\bibitem[\citeproctext]{ref-strobl2007bias}
Strobl, Carolin, Anne-Laure Boulesteix, and Thomas Augustin. 2007.
{``Unbiased Split Selection for Classification Trees Based on the {G}ini
Index.''} \emph{Computational Statistics \& Data Analysis} 52 (1):
483--501. \url{https://doi.org/10.1016/j.csda.2006.12.030}.

\bibitem[\citeproctext]{ref-turney1995stability}
Turney, Peter D. 1995. {``Bias and the Quantification of Stability.''}
\emph{Machine Learning} 20 (1): 23--33.
\url{https://doi.org/10.1007/BF00993473}.

\bibitem[\citeproctext]{ref-wang2017tsallis}
Wang, Yisen, Shu-Tao Xia, and Jia Wu. 2017. {``A Less-Greedy Two-Term
Tsallis Entropy Information Metric Approach for Decision Tree
Classification.''} \emph{Knowledge-Based Systems} 120: 34--42.
\url{https://doi.org/10.1016/j.knosys.2016.12.021}.

\bibitem[\citeproctext]{ref-yang2019pairgain}
Yang, Bo-Bo, Wei Gao, and Min Li. 2019. {``On the Robust Splitting
Criterion of Random Forest.''} \emph{2019 IEEE 19th International
Conference on Data Mining (ICDM)}, 1420--25.
\url{https://doi.org/10.1109/ICDM.2019.00184}.

\end{CSLReferences}

\end{document}
